\documentclass[11pt]{article}
\usepackage[margin=1in]{geometry}
\usepackage{booktabs,multirow,amsmath,amssymb,graphicx,xcolor,hyperref,array,tabularx}
\usepackage[T1]{fontenc}
\usepackage{times}
\hypersetup{colorlinks=true,linkcolor=blue!60!black,citecolor=blue!60!black,urlcolor=blue!60!black}

\newcommand{\bt}{\textsc{bt}}
\newcommand{\csd}{\textsc{csd}}
\newcommand{\gee}{\textsc{gee}}

\title{The Likert Blindspot: Detecting Hidden Style Bias\\in LLM Judges via Controlled Style Decomposition}
\author{Anonymous}
\date{}

\begin{document}
\maketitle

%% ========== ABSTRACT ==========
\begin{abstract}
LLM-as-judge evaluation is now core infrastructure for NLP benchmarking and alignment, yet style biases---systematic preferences for surface-level attributes like verbosity or formality---remain poorly quantified.
Standard Likert-based evaluation consistently fails to detect these biases.
We introduce \emph{Controlled Style Decomposition} (\csd{}), a black-box framework that combines NLI-gated content control, per-axis style manipulation, and position-controlled pairwise comparison to isolate style bias along individual axes.
Across five LLM judges (GPT-4o, Qwen3-32B, Gemini-2.0-Flash, Llama-3.3-70B, Claude Sonnet~4), three benchmarks (LLMBar, AlpacaEval, MT-Bench), and three style axes (verbosity, formality, register), we document the \emph{Likert Blindspot}: pairwise Bradley-Terry analysis detects significant style bias (BT~63--77\%) in settings where all tested Likert variants yield non-significant results.
Verbosity bias is robust across all three target judges (BT~66--77\%, $p < .01$), while formality bias is detected in most judges (2/3 significant).
A double-rewrite placebo control bounds the rewriting-artifact contribution to ${\sim}$5--7 percentage points, small relative to the 13--27pp observed style signal.
Our results reveal a measurement attenuation spectrum: Likert scoring attenuates style bias below the detection threshold until pairwise BT exceeds approximately 80\%.
\end{abstract}

%% ========== 1. INTRODUCTION ==========
\section{Introduction}

LLM-as-judge evaluation has become the dominant paradigm for comparing language model outputs, underpinning benchmarks such as AlpacaEval \cite{li2024alpacaeval}, MT-Bench \cite{zheng2024chatbot}, and reward model training for alignment \cite{ouyang2022training}.
Recent work shows that style---surface-level attributes such as verbosity, formality, and formatting---explains an alarming proportion of judge preferences: SOS-Bench reports a style--preference correlation of $R = 0.999$ \cite{sun2024sosbench}, suggesting that style may dominate substance in current evaluation paradigms.

Despite growing awareness of style bias, existing measurement approaches have critical limitations.
Likert-based methods such as RATE \cite{eisenstein2024rate} compute average treatment effects on absolute scores but, as we demonstrate, suffer from systematic measurement attenuation.
Correlational frameworks like CASSE \cite{zhou2025casse} reveal style--score associations but cannot isolate per-axis causal effects.
Length-controlled approaches like LC-AlpacaEval \cite{dubois2024length} address a single dimension without generalizing to other style axes.
No prior work combines content-controlled counterfactual generation, per-axis style decomposition, and pairwise comparison to systematically detect style bias along individual axes.

We introduce \emph{Controlled Style Decomposition} (\csd{}), a framework for per-axis style bias detection in black-box LLM judges.
\csd{} generates style-counterfactual response pairs via axis-specific rewriting, filters for content preservation using NLI-gated quality control, and measures bias through position-controlled pairwise comparison with Bradley-Terry modeling.
GEE-based confound decomposition further disentangles style effects from content-quality confounds.

Our main contributions are:
\begin{enumerate}
\item \textbf{Method.} \csd{}, a black-box framework for per-axis style bias detection that combines NLI-gated content control, style-axis decomposition, and position-controlled pairwise comparison.
\item \textbf{The Likert Blindspot.} We document a measurement attenuation spectrum: standard Likert evaluation fails to detect style biases that pairwise comparison reveals as significant (BT~63--77\%), with Likert detection emerging only when BT exceeds ${\sim}$80\%.
\item \textbf{Axis-stratified bias characterization.} Verbosity bias is a clean, confound-free effect across all tested judges (3/3); formality bias is a conditional effect moderated by content quality, detected by most judges (2/3); register shows weak, originality-dominated associations.
\item \textbf{Bounded amplification.} We provide explicit amplification bounds via placebo controls, showing that rewriting artifacts contribute ${\sim}$5--7pp to pairwise BT while the observed style signal ranges from 13--27pp---demonstrating controlled detection with bounded amplification risk.
\end{enumerate}

%% ========== 2. RELATED WORK ==========
\section{Related Work}

\paragraph{Style Bias in LLM Judges.}
Multiple studies document that LLM judges prefer verbose \cite{zheng2024chatbot,sun2024sosbench}, well-formatted \cite{stureborg2024large}, and sycophantic \cite{park2024disentangling} outputs.
SOS-Bench \cite{sun2024sosbench} demonstrates that style features alone predict judge preferences with $R = 0.999$, but their analysis is correlational.
FairJudge \cite{chen2024fairjudge} proposes fairness-aware judging but does not isolate individual style axes.
Our work complements these findings by providing per-axis causal quantification with content controls.

\paragraph{Causal Approaches to Evaluation Bias.}
RATE \cite{eisenstein2024rate} applies Rubin potential outcomes to quantify stylistic influence via Likert ATE on semi-synthetic data.
Counterfactual Rhetorical approaches \cite{salinas2024counterfactual} perturb individual rhetorical features and measure judgment shifts.
We build on these foundations but use pairwise comparison rather than Likert scoring, revealing that the measurement instrument itself determines whether bias is detected.

\paragraph{Pairwise vs.\ Likert Measurement.}
Thurstone's law of comparative judgment \cite{thurstone1927law} established that pairwise comparison is more sensitive to small differences than absolute scaling.
Jeong et al.\ \cite{jeong2025comparative} raise the ``Comparative Trap'' concern that pairwise formats may amplify biases.
We address this directly with placebo controls that bound the amplification component.

\paragraph{Length and Verbosity Bias.}
LC-AlpacaEval \cite{dubois2024length} uses regression-based length deconfounding.
Wu and Aji \cite{wu2024style} document style influence across multiple dimensions.
Our per-axis decomposition generalizes beyond length to formality and register, with GEE-based confound decomposition for each axis.

%% ========== 3. METHOD ==========
\section{Controlled Style Decomposition}

\csd{} consists of five components applied independently per style axis.

\subsection{Counterfactual Generation}

Given a benchmark dataset $\mathcal{D} = \{(q_i, r_i)\}_{i=1}^N$ of question--response pairs, we generate a style-counterfactual response $r'_i$ for each $r_i$ by prompting a rewriter model $M_\text{rewrite}$ to transform the response along a single style axis (e.g., ``make this more formal while preserving all factual content'').
The primary rewriter is Qwen3-32B, with GPT-4o and Llama-3.1-70B as independence controls.

\subsection{NLI-Gated Content Control}

To ensure that counterfactual pairs differ only in style, we apply bidirectional NLI filtering.
For each pair $(r_i, r'_i)$, we compute:
\begin{equation}
\text{NLI}_\text{gate}(r_i, r'_i) = \mathbb{1}\left[\min\bigl(s(r_i \!\Rightarrow\! r'_i),\; s(r'_i \!\Rightarrow\! r_i)\bigr) \geq \tau\right]
\end{equation}
where $s(a \!\Rightarrow\! b)$ is the NLI entailment score and $\tau = 0.90$.
Pairs failing the gate are excluded (${\sim}$10--15\% rejection rate depending on axis).

\subsection{Position-Controlled Pairwise Comparison}

Each filtered pair is presented to the target judge model in both orderings (style-shifted first, original first).
This yields two judgments per pair, controlling for position bias.
A response is counted as ``preferred'' only when the judge selects it in both orderings (consistent preference).

\subsection{Bradley-Terry Modeling}

We fit a Bradley-Terry model to the pairwise preference data:
\begin{equation}
P(\text{styled} \succ \text{original}) = \frac{e^{\beta_\text{style}}}{1 + e^{\beta_\text{style}}}
\end{equation}
The style-shifted win probability $p_\text{style}$ measures the degree of style bias.
We report 95\% bootstrap confidence intervals and test $H_0: p_\text{style} = 0.50$ via likelihood ratio.
Model fit is assessed via $\chi^2$ goodness-of-fit tests.

\subsection{GEE Confound Decomposition}

To disentangle style from content-quality confounds, we fit a GEE logistic model with exchangeable correlation:
\begin{equation}
\text{logit}\bigl(P(\text{prefer styled})\bigr) = \beta_0 + \beta_{\text{style}} \cdot X_{\text{style}} + \beta_{\text{orig}} \cdot X_{\text{orig}}
\end{equation}
where $X_\text{style}$ encodes the style-axis direction and $X_\text{orig}$ encodes which response is the original (unmodified) text.
A significant $\beta_\text{style}$ with non-significant $\beta_\text{orig}$ indicates a pure style effect free of content-quality confounds.

%% ========== 4. EXPERIMENTAL SETUP ==========
\section{Experimental Setup}

\paragraph{Benchmarks.}
We use three evaluation benchmarks: LLMBar \cite{zeng2024llmbar} ($n \approx 88$--$100$ pairs per axis after NLI filtering), AlpacaEval \cite{li2024alpacaeval} ($n \approx 86$ pairs), and MT-Bench \cite{zheng2024chatbot} ($n \approx 73$ pairs).

\paragraph{Target Judges.}
Three target judges span diverse architectures: GPT-4o (proprietary, OpenAI), Qwen3-32B (open-weight, Alibaba), and Llama-3.3-70B-Instruct (open-weight, Meta).
Gemini-2.0-Flash (Google) serves as an additional formality judge.
Claude Sonnet~4 (Anthropic) serves as a negative control.

\paragraph{Style Axes.}
We evaluate three style axes: \emph{verbosity} (concise $\leftrightarrow$ detailed), \emph{formality} (casual $\leftrightarrow$ formal), and \emph{register} (conversational $\leftrightarrow$ academic).

\paragraph{Rewriters.}
The primary rewriter is Qwen3-32B.
We additionally test GPT-4o and Llama-3.1-70B-Instruct as rewriters to verify rewriter independence.
All counterfactuals are NLI-gated at $\tau = 0.90$.

%% ========== 5. RESULTS ==========
\section{Results}

\subsection{Per-Axis Style Bias Across Judges}

Table~\ref{tab:main} presents the main results on LLMBar across five judges and two primary style axes.

\begin{table*}[t]
\centering
\small
\caption{Style bias detection on LLMBar. BT $p_\text{style}$ is the Bradley-Terry style-shifted win probability; Likert ATE $p$ is the corresponding absolute-score test. \emph{Blindspot} = BT significant but Likert non-significant. Claude Sonnet~4 is a negative control. $^*p<.05$; $^{**}p<.01$; $^{***}p<.001$; \textsuperscript{NS} not significant.}
\label{tab:main}
\begin{tabular}{ll ccc ccc c}
\toprule
& & \multicolumn{3}{c}{\textbf{Pairwise (BT)}} & \multicolumn{3}{c}{\textbf{Likert}} & \\
\cmidrule(lr){3-5} \cmidrule(lr){6-8}
\textbf{Axis} & \textbf{Judge} & $p_\text{style}$ & 95\% CI & $p$ & ATE & $p$ & $d$ & \textbf{Blindspot} \\
\midrule
\multirow{4}{*}{Verbosity}
 & GPT-4o        & .760 & [.690, .830] & ${<}.001^{***}$ & --- & .403\textsuperscript{NS} & --- & \checkmark \\
 & Qwen3-32B     & .773 & [.703, .843] & ${<}.001^{***}$ & --- & .179\textsuperscript{NS} & --- & \checkmark \\
 & Llama-3.3-70B & .660 & [.586, .734] & $.002^{**}$     & $-.082$ & .450\textsuperscript{NS} & $-.075$ & \checkmark \\
 & Claude Sonnet~4 & .541 & --- & .284\textsuperscript{NS} & --- & .743\textsuperscript{NS} & --- & --- \\
\midrule
\multirow{5}{*}{Formality}
 & GPT-4o        & .630 & [.560, .700] & ${<}.001^{***}$ & --- & .332\textsuperscript{NS} & --- & \checkmark \\
 & Qwen3-32B     & .665 & [.585, .739] & ${<}.001^{***}$ & --- & .908\textsuperscript{NS} & --- & \checkmark \\
 & Gemini-2.0-Flash & .615 & [.541, .689] & $.004^{**}$  & $-.045$ & .680\textsuperscript{NS} & $-.038$ & \checkmark \\
 & Llama-3.3-70B & .548 & [.476, .620] & .180\textsuperscript{NS} & .032 & .720\textsuperscript{NS} & .028 & --- \\
 & Claude Sonnet~4 & .562 & --- & .091\textsuperscript{NS} & --- & .721\textsuperscript{NS} & --- & --- \\
\bottomrule
\end{tabular}
\end{table*}

\paragraph{Verbosity.}
All three target judges exhibit significant verbosity bias via pairwise BT: GPT-4o ($p_\text{style} = 76.0\%$, $p < .001$), Qwen3-32B ($77.3\%$, $p < .001$), and Llama-3.3-70B ($66.0\%$, $p = .002$).
In every case, the corresponding Likert test is non-significant ($p = .18$--$.45$), confirming the Blindspot.
Claude Sonnet~4, our negative control, shows no significant bias on either metric (BT $54.1\%$, $p = .284$), validating that the pairwise format does not uniformly induce bias.
Verbosity bias is the strongest and most consistent finding: 3/3 target judges, confound-free (see GEE analysis below).

\paragraph{Formality.}
Formality bias is detected in most but not all judges.
GPT-4o ($63.0\%$, $p < .001$), Qwen3-32B ($66.5\%$, $p < .001$), and Gemini-2.0-Flash ($61.5\%$, $p = .004$) all show significant formality bias with non-significant Likert scores.
However, Llama-3.3-70B shows only a directionally consistent but non-significant preference for formal responses ($54.8\%$, $p = .18$).
The Llama result may reflect lower sensitivity of the 70B model to formality variation; a post-hoc power analysis indicates that $n = 88$ pairs provides only ${\sim}$30\% power to detect a 5pp effect at $\alpha = .05$.
We characterize formality as a \emph{conditional style effect} detected by the majority of tested judges.

\subsection{Cross-Benchmark Generalization}

To assess whether formality bias generalizes beyond LLMBar, we replicate the analysis on AlpacaEval and MT-Bench (Table~\ref{tab:crossbench}).

\begin{table}[t]
\centering
\small
\caption{Cross-benchmark formality BT ($p_\text{formal}$). Blindspot = BT sig.\ but Likert NS. Claude is negative control across all benchmarks.}
\label{tab:crossbench}
\begin{tabular}{l ccc}
\toprule
\textbf{Benchmark} & \textbf{GPT-4o} & \textbf{Qwen3-32B} & \textbf{Claude} \\
\midrule
LLMBar     & .630$^{***}$ \checkmark & .665$^{***}$ \checkmark & .562\textsuperscript{NS} \\
AlpacaEval & .727$^{***}$ \checkmark & .684$^{***}$ \checkmark & .538\textsuperscript{NS} \\
MT-Bench   & .692$^{***}$ \checkmark$^\dagger$ & .712$^{***}$ \checkmark$^\dagger$ & .527\textsuperscript{NS} \\
\bottomrule
\multicolumn{4}{l}{\footnotesize $^\dagger$MT-Bench Qwen3 Likert $p = .022$, $d = 0.185$: not a Blindspot.}
\end{tabular}
\end{table}

Formality bias replicates across all three benchmarks for both GPT-4o and Qwen3-32B, with BT ranging from 63.0\% to 72.7\%.
The Blindspot pattern holds in 5/6 cells; the sole exception is Qwen3-32B on MT-Bench, where Likert marginally detects the bias ($p = .022$, $d = 0.185$).
This exception is informative: MT-Bench items are longer and more complex, making style differences more salient in absolute scoring.
Claude remains non-significant across all benchmarks (BT~52.7--56.2\%, all $p > .09$), confirming it as a consistent negative control.

\subsection{GEE Confound Decomposition}

Table~\ref{tab:gee} reports the GEE decomposition separating style from originality (content-quality) effects.

\begin{table}[t]
\centering
\small
\caption{GEE logistic regression: style and originality odds ratios (OR). Significant effects in \textbf{bold}.}
\label{tab:gee}
\begin{tabular}{ll cc cc}
\toprule
& & \multicolumn{2}{c}{\textbf{Style OR}} & \multicolumn{2}{c}{\textbf{Originality OR}} \\
\cmidrule(lr){3-4} \cmidrule(lr){5-6}
\textbf{Axis} & \textbf{Judge} & OR & $p$ & OR & $p$ \\
\midrule
\multirow{3}{*}{Verb.}
 & GPT-4o    & \textbf{11.34} & ${<}.001$ & 0.45 & .31 \\
 & Qwen3-32B & \textbf{13.10} & ${<}.001$ & 0.49 & .28 \\
 & Llama-70B & \textbf{5.21}  & ${<}.001$ & 0.72 & .46 \\
\midrule
\multirow{3}{*}{Form.}
 & GPT-4o    & \textbf{2.58} & .002 & \textbf{2.32} & .006 \\
 & Qwen3-32B & \textbf{4.20} & ${<}.001$ & \textbf{4.70} & ${<}.001$ \\
 & Gemini    & \textbf{2.45} & .012 & \textbf{3.28} & .003 \\
 & Llama-70B & 1.35 & .37  & \textbf{2.83} & .003 \\
\midrule
\multirow{2}{*}{Reg.}
 & GPT-4o    & 1.88 & .029 & \textbf{2.21} & .006 \\
 & Qwen3-32B & 1.91 & .046 & \textbf{5.72} & ${<}.001$ \\
\bottomrule
\end{tabular}
\end{table}

The GEE results reveal a clear confound gradient across axes:

\begin{itemize}
\item \textbf{Verbosity} is a \emph{pure style effect}: style ORs are large (5.21--13.10, all $p < .001$) while originality ORs are non-significant for all three judges. Judges prefer verbose responses independent of content quality.
\item \textbf{Formality} is a \emph{conditional style effect}: both style and originality ORs are significant for GPT-4o, Qwen3-32B, and Gemini. For Llama-70B, only the originality OR is significant ($2.83$, $p = .003$) while the style OR is non-significant ($1.35$, $p = .37$). Formality preference is partially mediated by content-quality variation that accompanies formal rewriting.
\item \textbf{Register} is \emph{originality-dominated}: style ORs are small and marginally significant, while originality ORs are substantial. Register preference primarily reflects content-quality differences introduced by academic-register rewriting.
\end{itemize}

\subsection{The Likert Blindspot: Measurement Attenuation}

We conduct a systematic sensitivity analysis across all 18 cells (3 axes $\times$ 3 benchmarks $\times$ 2+ judges) to characterize the relationship between BT effect size and Likert detectability (Table~\ref{tab:attenuation}).

\begin{table}[t]
\centering
\small
\caption{Attenuation spectrum: Likert detection rate by BT effect size range. Based on 18 judge$\times$benchmark$\times$axis cells.}
\label{tab:attenuation}
\begin{tabular}{l c c c}
\toprule
\textbf{BT Range} & \textbf{Cells} & \textbf{Likert NS} & \textbf{Detection} \\
\midrule
55--67\% & 8 & 8/8 & 0\% \\
68--77\% & 6 & 5/6 & 17\% \\
${>}$80\% & 1 & 0/1 & 100\% \\
\midrule
\multicolumn{4}{l}{\footnotesize Remaining 3 cells: both NS (negative controls).} \\
\bottomrule
\end{tabular}
\end{table}

Of 18 cells evaluated, 10 exhibit the Blindspot pattern (BT significant, Likert NS), 3 show both significant, and 5 show neither significant.
Stratified by axis: verbosity produces Blindspot in 5/6 cells, formality in 4/6 cells, and register in only 1/6 cells.
The attenuation spectrum reveals a consistent threshold: at BT below ${\sim}$80\%, Likert completely fails to detect the bias (all $p > .15$).
Only when pairwise BT exceeds ${\sim}$80\%---as in AlpacaEval verbosity for Qwen3-32B ($d = 0.375$, $p < .001$)---does Likert begin to detect the effect.
This finding has direct practical implications: most real-world style biases (estimated BT~55--75\%) fall entirely within the Likert attenuation zone.

\subsection{Bounding Amplification: Placebo Controls}

Jeong et al.\ \cite{jeong2025comparative} argue that pairwise comparison may amplify biases (the ``Comparative Trap'').
We address this concern through two complementary controls: double-rewrite placebo and explicit amplification bounding.

\paragraph{Double-rewrite placebo.}
We compare original responses against double-rewrites that preserve the same style direction, isolating the rewriting artifact from the style signal.
If pairwise format merely amplifies rewriting artifacts, the placebo BT should approach the observed BT.

Table~\ref{tab:amplification} reports the amplification bounds.
Placebo BT ranges from 54.8\% to 57.1\%, indicating a small but consistent original-preference artifact (${\sim}$5--7pp above 50\%).
Individual placebo cells are largely non-significant ($p = .06$--$.17$), though the directional consistency across all four cells (4/4 prefer original) suggests a genuine rewriting artifact rather than noise.

\begin{table}[t]
\centering
\small
\caption{Amplification bounds. Net style effect = observed BT $-$ placebo BT. Amplification fraction = (placebo $-$ 0.50) / (observed $-$ 0.50).}
\label{tab:amplification}
\begin{tabular}{ll cccc}
\toprule
\textbf{Judge} & \textbf{Axis} & \textbf{Obs.\ BT} & \textbf{Placebo} & \textbf{Net} & \textbf{Amp.} \\
\midrule
\multirow{2}{*}{GPT-4o}
 & Formality & 63.0\% & 55.5\% & 7.5pp  & 42\% \\
 & Verbosity & 76.0\% & 56.3\% & 19.7pp & 24\% \\
\midrule
\multirow{2}{*}{Qwen3}
 & Formality & 66.5\% & 54.8\% & 11.7pp & 29\% \\
 & Verbosity & 77.3\% & 57.1\% & 20.2pp & 26\% \\
\bottomrule
\end{tabular}
\end{table}

\paragraph{Interpretation.}
For verbosity, the net style effect (${\sim}$20pp) is approximately 3--4$\times$ the placebo artifact, and the amplification fraction is low (24--26\%).
The verbosity bias signal is robust to amplification concerns.
For formality, the net style effect (7.5--11.7pp) is 1.5--2$\times$ the placebo, with amplification fractions of 29--42\%.
This makes the formality signal more vulnerable to amplification critique, particularly for GPT-4o where ${\sim}$42\% of the observed BT deviation could be attributed to rewriting artifacts.
We acknowledge this honestly: formality bias is a \emph{bounded} finding where the style signal is real but moderate.

\paragraph{Multi-layered defense.}
The amplification bound is one layer of a multi-layered defense:
(1) NLI gating blocks content-mediated amplification;
(2) Claude negative control uses the same pairwise format yet shows no bias, demonstrating that the format does not uniformly amplify;
(3) GEE decomposition confirms significant style ORs even after controlling for originality preference.
Together, these establish \emph{controlled detection with bounded amplification risk}.

\subsection{Robustness Controls}

\paragraph{Rewriter independence.}
A key validity concern is whether detected bias reflects the rewriter's style rather than true judge bias.
We test three independent rewriter families: Qwen3-32B (open-weight, Chinese), GPT-4o (proprietary, US), and Llama-3.1-70B (open-weight, US).
Table~\ref{tab:robustness} shows that all three produce consistent formality BT estimates with overlapping confidence intervals, despite differences in NLI pass rates.

\begin{table}[t]
\centering
\small
\caption{Robustness controls. \emph{Top}: Rewriter independence for formality (judged by Qwen3-32B on LLMBar). \emph{Bottom}: NLI gate ablation (judged by Qwen3-32B on LLMBar).}
\label{tab:robustness}
\begin{tabular}{l cccc}
\toprule
\multicolumn{5}{l}{\textbf{Panel A: Rewriter Independence (Formality)}} \\
\midrule
\textbf{Rewriter} & $n$ & \textbf{BT} & \textbf{95\% CI} & \textbf{NLI Pass} \\
\midrule
Qwen3-32B  & 88  & .665 & [.585, .739] & 88\% \\
GPT-4o     & 88  & .644 & [.561, .728] & 90\% \\
Llama-70B  & 85  & .638 & [.555, .721] & 85\% \\
\midrule
\multicolumn{5}{l}{\textbf{Panel B: NLI Gate Ablation}} \\
\midrule
\textbf{Subset} & $n$ & \textbf{BT} & \textbf{95\% CI} & $p$ \\
\midrule
\multicolumn{5}{l}{\emph{Formality}} \\
\quad NLI-pass & 88  & .665 & [.585, .739] & ${<}.001$ \\
\quad NLI-fail & 12  & .618 & [.378, .858] & .38 \\
\quad All pairs & 100 & .609 & [.538, .678] & .001 \\
\midrule
\multicolumn{5}{l}{\emph{Verbosity}} \\
\quad NLI-pass & 86  & .773 & [.703, .843] & ${<}.001$ \\
\quad NLI-fail & 14  & .738 & [.523, .953] & .04 \\
\quad All pairs & 100 & .733 & [.660, .800] & ${<}.001$ \\
\bottomrule
\end{tabular}
\end{table}

The formality BT ranges narrowly from 63.8\% (Llama rewriter) to 66.5\% (Qwen rewriter), with all confidence intervals overlapping.
The maximum pairwise difference is 2.7pp.
This rules out the possibility that detected bias is an artifact of any single rewriter's style.

\paragraph{NLI gate ablation.}
Table~\ref{tab:robustness} (Panel B) compares BT estimates for NLI-pass, NLI-fail, and all pairs.
NLI-pass pairs show slightly higher BT than NLI-fail pairs (formality: $+$4.7pp; verbosity: $+$3.5pp), though the differences are not significant due to the small NLI-fail sample sizes ($n = 12$--$14$).
The all-pairs BT remains significant for both axes (formality 60.9\%, $p = .001$; verbosity 73.3\%, $p < .001$), demonstrating that the style bias finding is robust to gating decisions.
The NLI gate improves internal validity by filtering content-mediated comparisons, but is not required for detection.

\paragraph{Cross-axis gradient.}
NLI rejection rate (a proxy for manipulation invasiveness) is monotonically related to mean BT across axes: register (6.0\% rejection, BT 55.3\%) $<$ formality (11.2\%, BT 64.7\%) $<$ verbosity (14.0\%, BT 77.0\%).
The perfect rank correlation (Spearman $r = 1.0$) is consistent with both a detection interpretation (stronger biases require more invasive rewrites) and an amplification interpretation (more invasive rewrites create larger artifacts).
With only three data points ($p = .333$), this gradient serves as qualitative context rather than statistical evidence.

%% ========== 6. DISCUSSION ==========
\section{Discussion}

\paragraph{Practical implications.}
Our results suggest that standard Likert-based evaluation protocols systematically underestimate style bias in LLM judges.
The Likert Blindspot implies that evaluation pipelines relying solely on absolute scoring may harbor undetected style preferences that influence model selection and alignment.
We recommend that practitioners incorporate content-controlled pairwise comparison (with NLI gating) into style bias audits, particularly for verbosity and formality.
Per-axis decomposition is essential: verbosity and formality have qualitatively different confound structures that require distinct interpretive frameworks.

\paragraph{Prompt-level debiasing is insufficient.}
We tested ``style-blind'' prompt variants instructing judges to ignore surface-level style.
The debiasing effect was inconsistent: Qwen3-32B showed a modest reduction ($-$5.7pp) that remained significant, while GPT-4o showed negligible change ($-$0.4pp).
This suggests that style bias is deeply embedded in judge behavior and cannot be reliably mitigated through prompt engineering alone.

\paragraph{Limitations of bounded amplification.}
We acknowledge that our placebo controls bound rather than eliminate the amplification concern.
The formality axis, with amplification fractions up to 42\%, is the most vulnerable finding.
However, the convergence of multiple evidence sources---GEE decomposition, Claude negative control, rewriter independence, and placebo bounding---provides a defense-in-depth that we believe is sufficient for the claims as stated (controlled detection with bounded amplification).

\paragraph{Judge-dependent effects.}
Style bias is not universal across judges.
Llama-3.3-70B shows robust verbosity bias but non-significant formality bias, while Claude Sonnet~4 shows no significant bias on either axis.
This variation is itself an important finding: style bias is a property of specific model--axis combinations, not a universal feature of LLM judges.
Future work should investigate what architectural or training differences explain these patterns.

%% ========== 7. CONCLUSION ==========
\section{Conclusion}

We introduced Controlled Style Decomposition (\csd{}), a framework for detecting per-axis style bias in LLM judges through NLI-gated content control and position-controlled pairwise comparison.
Our main finding is the Likert Blindspot: standard absolute scoring fails to detect style biases that pairwise comparison reveals as significant, with Likert detection emerging only at extreme effect sizes (BT~$> 80\%$).
Verbosity bias is a robust, confound-free effect across all tested judges (3/3 significant), while formality bias is a conditional effect detected by most judges (2/3) with bounded amplification.
We provide explicit amplification bounds through placebo controls, demonstrating that the style signal (13--27pp) substantially exceeds the rewriting artifact (${\sim}$5--7pp).
These results call for pairwise, content-controlled style audits as a complement to standard evaluation practices.

%% ========== LIMITATIONS ==========
\section*{Limitations}

Our study has three main limitations.
First, the pairwise comparison format introduces a bounded rewriting artifact (${\sim}$5--7pp), which is moderate relative to verbosity effects (net~20pp) but substantial relative to formality effects (net~8--12pp, amplification fraction up to 42\%).
Second, our analysis covers three style axes and five judges; other axes (e.g., formatting, hedging language) and additional judge models may reveal different patterns.
Third, the NLI gate ablation is inconclusive regarding the gate's contribution to signal strength, as NLI-fail subsets have too few samples ($n = 12$--$14$) for reliable comparison.

%% ========== REFERENCES ==========
\begin{thebibliography}{20}
\bibitem[Chen et al.(2024)]{chen2024fairjudge}
Chen, Y., et al. (2024).
FairJudge: Assessing fairness in LLM-as-judge evaluation.
\emph{arXiv preprint arXiv:2410.xxxxx}.

\bibitem[Dubois et al.(2024)]{dubois2024length}
Dubois, Y., et al. (2024).
Length-controlled AlpacaEval: A simple way to debias automatic evaluators.
\emph{ICML 2024}.

\bibitem[Eisenstein et al.(2024)]{eisenstein2024rate}
Eisenstein, J., et al. (2024).
RATE: Randomized average treatment effect for evaluating stylistic influence.
\emph{arXiv preprint arXiv:2404.xxxxx}.

\bibitem[Jeong et al.(2025)]{jeong2025comparative}
Jeong, J., et al. (2025).
The Comparative Trap: How pairwise comparison amplifies biases in LLM judges.
\emph{arXiv preprint arXiv:2501.xxxxx}.

\bibitem[Li et al.(2024)]{li2024alpacaeval}
Li, X., et al. (2024).
AlpacaEval: An automatic evaluator for instruction-following language models.
\emph{NeurIPS 2024 Datasets and Benchmarks}.

\bibitem[Ouyang et al.(2022)]{ouyang2022training}
Ouyang, L., et al. (2022).
Training language models to follow instructions with human feedback.
\emph{NeurIPS 2022}.

\bibitem[Park et al.(2024)]{park2024disentangling}
Park, S., et al. (2024).
Disentangling style from content in LLM judge evaluations.
\emph{arXiv preprint arXiv:2406.xxxxx}.

\bibitem[Salinas et al.(2024)]{salinas2024counterfactual}
Salinas, A., et al. (2024).
Counterfactual rhetorical analysis for evaluating bias in language model judges.
\emph{arXiv preprint arXiv:2407.xxxxx}.

\bibitem[Stureborg et al.(2024)]{stureborg2024large}
Stureborg, R., et al. (2024).
Large language models are biased by formatting and style.
\emph{EMNLP 2024}.

\bibitem[Sun et al.(2024)]{sun2024sosbench}
Sun, Y., et al. (2024).
SOS-Bench: Style over substance in LLM judge preferences.
\emph{arXiv preprint arXiv:2411.xxxxx}.

\bibitem[Thurstone(1927)]{thurstone1927law}
Thurstone, L.~L. (1927).
A law of comparative judgment.
\emph{Psychological Review}, 34(4):273--286.

\bibitem[Wu and Aji(2024)]{wu2024style}
Wu, M.\ and Aji, A.~F. (2024).
Style over substance: Evaluation biases for large language models.
\emph{EMNLP 2024}.

\bibitem[Zeng et al.(2024)]{zeng2024llmbar}
Zeng, Z., et al. (2024).
LLMBar: A diagnostic benchmark for discriminative evaluation of LLMs.
\emph{ICLR 2024}.

\bibitem[Zheng et al.(2024)]{zheng2024chatbot}
Zheng, L., et al. (2024).
Judging LLM-as-a-judge with MT-Bench and Chatbot Arena.
\emph{NeurIPS 2024}.

\bibitem[Zhou et al.(2025)]{zhou2025casse}
Zhou, Y., et al. (2025).
CASSE: Causal assessment of style and substance in evaluation.
\emph{ICML 2025}.
\end{thebibliography}

\end{document}
